% ======================================================================
\section{States, Hamiltonian and light-cone matrix elements}\label{app:hamiltonian}
% ======================================================================

In this appendix we collect the definitions of the states, the parts of the eikonal light-cone Hamiltonian that are needed, and the matrix elements used in the LCPT calculations of Appendix~\ref{app:wavefunctions}. They are those of~\cite{Radhakrishnan:2026soft,Lublinsky:2016meo}, restricted to pure Yang--Mills theory.

\subsection{States and normalisation}\label{app:states}

The soft vacuum $\ket0$ is annihilated by all soft annihilation operators, $a^a_i(k^+,\v k)\ket0=0$ for $\Lambda<k^+<\vee$, and its energy is set to zero, $E_0=0$. The one- and two-gluon states are
\begin{align}
\ket{g^a_i(k)}&=\frac{a^{\dagger a}_i(k)}{(2\pi)^{3/2}}\ket0\,,&
\ket{g^a_i(k^+,\v z)}&=\frac{a^{\dagger a}_i(k^+,\v z)}{(2\pi)^{1/2}}\ket0\,,\notag\\
\ket{g^a_i(k)\,g^b_j(p)}&=\frac{a^{\dagger a}_i(k)\,a^{\dagger b}_j(p)}{(2\pi)^3}\ket0\,,&
\ket{g^a_i(k^+,\v z)\,g^b_j(p^+,\v z')}&=\frac{a^{\dagger a}_i(k^+,\v z)\,a^{\dagger b}_j(p^+,\v z')}{2\pi}\ket0\,,
\end{align}
with energies $E_g(k)=\v k^2/2k^+$ and $E_{gg}(k,p)=\v k^2/2k^++\v p^2/2p^+$, and normalisation
\begin{align}
\braket{g^b_j(p)|g^a_i(k)}&=\delta^{ab}\delta_{ij}\,\delta(k^+-p^+)\,\delta^{(2)}(\v k-\v p)\,,\notag\\
\braket{g^b_j(p^+,\v z')|g^a_i(k^+,\v z)}&=\delta^{ab}\delta_{ij}\,\delta(k^+-p^+)\,\delta^{(2)}(\v z-\v z')\,.
\end{align}
The coordinate-space state differs from the Fourier transform of the momentum-space state by a factor $2\pi$, $\ket{g(k^+,\v z)}=2\pi\int\frac{d^2k}{(2\pi)^2}e^{-i\v k\cdot\v z}\ket{g(k)}$. Consequently $a^a_i(k^+,\v k)\ket{g^b_j(q)}=(2\pi)^{3/2}\delta^{ab}\delta_{ij}\delta^{(3)}(k-q)\ket0$.

\subsection{Soft and valence modes}

The gauge field in the light-cone gauge is decomposed into soft and valence components,
\begin{align}
\underline{A}^a_i(x)&=\int_\Lambda^\vee\frac{dk^+}{2\pi}\int\frac{d^2k}{(2\pi)^2}\frac{1}{\sqrt{2k^+}}\Big(a^a_i(k^+,\v k)e^{-ik\cdot x}+a^{\dagger a}_i(k^+,\v k)e^{ik\cdot x}\Big)\,,\notag\\
\overline{A}^a_i(x)&=\int_\vee^\infty\frac{dk^+}{2\pi}\int\frac{d^2k}{(2\pi)^2}\frac{1}{\sqrt{2k^+}}\Big(a^a_i(k^+,\v k)e^{-ik\cdot x}+a^{\dagger a}_i(k^+,\v k)e^{ik\cdot x}\Big)\,.
\end{align}
The valence charge density is the colour charge density of the valence modes~\cite{Radhakrishnan:2026soft},
\begin{align}
\rho^a(-\v p)=-if^{abc}\int_\vee^\infty\frac{dk^+}{2\pi}\int\frac{d^2k}{(2\pi)^2}\,a^{\dagger b}_j(k^+,\v k)\,a^c_j(k^+,\v k+\v p)\,,
\end{align}
and satisfies the algebra~\eqref{eq:rho}. Inserting the decomposition into the light-cone Hamiltonian and dropping all terms suppressed by powers of the ratio of soft to valence longitudinal momenta (in particular all inverse longitudinal derivatives acting on valence fields), one obtains the interaction terms~\cite{Lublinsky:2016meo,Radhakrishnan:2026soft}
\begin{align}
H_g&=-g\int dx^-d^2x\,\big(\partial_i\underline A^a_i\big)\,f^{abc}\,\frac{1}{\partial^+}\big(\overline A^b_j\partial^+\overline A^c_j\big)\,,\notag\\
H_{ggg}&=-gf^{abc}\int dx^-d^2x\,\Big[\big(\partial_i\underline A^a_j\big)\underline A^b_i\underline A^c_j+\big(\partial_i\underline A^a_i\big)\frac{1}{\partial^+}\big(\underline A^b_j\partial^+\underline A^c_j\big)\Big]\,,\notag\\
H_{gg\text{-inst}}&=g^2f^{abc}f^{ade}\int dx^-d^2x\,\frac{1}{\partial^+}\big(\underline A^b_i\partial^+\underline A^c_i\big)\,\frac{1}{\partial^+}\big(\overline A^d_j\partial^+\overline A^e_j\big)\,,\notag\\
H_{gggg}&=\frac{g^2}{4}f^{abc}f^{ade}\int dx^-d^2x\,\Big[\underline A^b_i\underline A^c_j\underline A^d_i\underline A^e_j-2\,\frac{1}{\partial^+}\big(\underline A^b_i\partial^+\underline A^c_i\big)\frac{1}{\partial^+}\big(\underline A^d_j\partial^+\underline A^e_j\big)\Big]\,.
\end{align}
In terms of the soft creation and annihilation operators $H_g$ takes the form~\eqref{eq:Hg}, and
\begin{align}
H_{gg\text{-inst}}=\ &\int_\Lambda^\vee\frac{dp^+}{2\pi}\frac{dq^+}{2\pi}\int\frac{d^2p}{(2\pi)^2}\frac{d^2q}{(2\pi)^2}\frac{ig^2f^{abc}}{\sqrt{p^+q^+}}\notag\\&\times\Bigg(\frac{q^+\rho^a(-\v p-\v q)}{2(p^++q^+)^2}\,a^{\dagger b}_i(p^+,\v p)\,a^{\dagger c}_i(q^+,\v q)-{\rm h.c.}
\notag\\&
-\frac{(p^++q^+)\,\rho^a(-\v p+\v q)}{2(p^+-q^+)^2}\,a^{\dagger b}_i(p^+,\v p)\,a^c_i(q^+,\v q)\Bigg)\,.
\end{align}
The four-gluon interaction $H_{gggg}$ does not enter the coefficients used in this paper.
% CHOICE FOR THE AUTHORS: confirm that H_gggg (normal ordered) gives no O(g^3) one-gluon contribution (no H_gggg diagram appears in the calculation files).

\subsection{Matrix elements}\label{app:matrixelements}

The matrix elements of the interaction terms between the free states, which enter the LCPT expansion~\eqref{eq:lcpt}, are~\cite{Radhakrishnan:2026soft}
\begin{align}
\braket{g^a_i(k)|H_g|0}&=\frac{g\,k^i\rho^a(-\v k)}{4\pi^{3/2}|k^+|^{3/2}}\,,\label{eq:m1}\\
\braket{g^c_l(q)\,g^b_j(p)|H_g|g^a_i(k)}&=\delta^{ac}\delta_{il}\delta^{(3)}(k-q)\frac{g\,p^j\rho^b(-\v p)}{4\pi^{3/2}|p^+|^{3/2}}\notag\\&\quad+\delta^{ab}\delta_{ij}\delta^{(3)}(k-p)\frac{g\,q^l\rho^c(-\v q)}{4\pi^{3/2}|q^+|^{3/2}}\,,\label{eq:m2}\\
\braket{g^c_l(q)\,g^b_j(p)|H_{ggg}|g^a_i(k)}&=\frac{igf^{abc}\,\delta^{(3)}(k-p-q)}{8\pi^{3/2}\sqrt{k^+p^+q^+}}\Bigg[\Big(p^i-q^i+\frac{q^+-p^+}{k^+}k^i\Big)\delta_{jl}\notag\\&\hspace{1cm}+\Big(k^j+q^j-\frac{k^++q^+}{p^+}p^j\Big)\delta_{il}
\notag\\&\hspace{1cm}
+\Big(\frac{k^++p^+}{q^+}q^l-p^l-k^l\Big)\delta_{ij}\Bigg]\,,\label{eq:m3}\\
\braket{g^c_j(q)\,g^b_i(p)|H_{gg\text{-inst}}|0}&
=\frac{ig^2f^{abc}\,(q^+-p^+)\,\delta_{ij}\,\rho^a(-\v p-\v q)}{2(2\pi)^3\sqrt{p^+q^+}\,(p^++q^+)^2}\,,\label{eq:m4}\\
\braket{g^c_j(q)|H_{gg\text{-inst}}|g^b_i(p)}&
=\frac{ig^2f^{abc}\,(p^++q^+)\,\delta_{ij}\,\rho^a(\v p-\v q)}{2(2\pi)^3\sqrt{p^+q^+}\,(p^+-q^+)^2}\,.\label{eq:m5}
\end{align}
Eq.~\eqref{eq:m1} is the emission of a soft gluon by the valence charges, Eq.~\eqref{eq:m2} the emission of a second gluon in the presence of a first one, Eq.~\eqref{eq:m3} the triple-gluon vertex among soft gluons, Eq.~\eqref{eq:m4} the instantaneous emission of two soft gluons by the valence charges and Eq.~\eqref{eq:m5} the instantaneous interaction of a soft gluon with the valence charges, in which the longitudinal momentum $p^+-q^+$ is exchanged with the valence modes. All matrix elements are operators on the valence Hilbert space.

\subsection{Rules of light-cone perturbation theory}\label{app:lcptrules}

For the dressing of the soft vacuum the energy denominators are $E_0-E_n=-E_n$, Eq.~\eqref{eq:lcpt}. For the dressing of a one-gluon state $\ket{g^a_i(k)}$ the energy of the initial state is $E_g(k)$ and the denominators are $E_g(k)-E_n$; for example, the instantaneous transition of the gluon into a gluon of momentum $p$ is
\begin{align}
\ket{g^a_i(k)}\to\int dp^+\!\int d^2p\;\ket{g^b_j(p)}\,\frac{\braket{g^b_j(p)|H_{gg\text{-inst}}|g^a_i(k)}}{E_g(k)-E_g(p)}\,.
\end{align}
These denominators vanish inside the integration region, on $p^+\v k^2=k^+\v p^2$ in momentum space or on $p^+(\v x-\v u)^2=k^+(\v x-\v z)^2$ in coordinate space; they are understood as principal values (Appendix~\ref{app:distributions}). The longitudinal momentum of every soft gluon, real or intermediate, is restricted to $\Lambda<p^+<\vee$, and the longitudinal momentum exchanged with the valence charges in Eqs.~\eqref{eq:m2} and~\eqref{eq:m5} satisfies $|p^+-q^+|>\Lambda$. The kinematic regions allowed for the transition $g(k^+)\to g(p^+)$ through the valence charges are shown in Figure~\ref{fig:kinem}.

\begin{figure}[t]
\centering
\tikzfig{kinem}
\caption{Allowed longitudinal momenta for the transition of a soft gluon of momentum $k^+$ into a soft gluon of momentum $p^+$ by the interaction with the valence charges, for $p^+>k^+$. The upper (purple) region corresponds to $\Lambda<p^+-k^+<\vee$, the green region to $\Lambda<k^+<\vee$ and $\Lambda<p^+<\vee$; the allowed region is their intersection. The case $p^+<k^+$, in which the gluon transfers longitudinal momentum to the valence charges, is the mirror image.}
\label{fig:kinem}
\end{figure}
